% !PRINTED 518
\weberpage{518}
3. \gesperrt{Ein ganzes Functional $\omega$ im Körper $\Omega$ ist nur auf eine Weise in Primfactoren zerlegbar, wenn associirte Primfactoren als nicht verschieden betrachtet werden. Associirte Functionale enthalten dieselben Primfactoren.}

Nehmen wir nämlich an, es seien die beiden Producte von Primfactoren
\weq{6}{\pi_1 \pi_2 \ldots \pi_\nu, \qquad \varkappa_1 \varkappa_2 \ldots \varkappa_\mu}
mit einander associirt, so ist das Product $\pi_1 \pi_2 \ldots \pi_\nu$ durch den Primfactor $\varkappa_1$ theilbar, und es muss daher, nach \S.\,140, 2., einer der Factoren, etwa $\pi_1$, durch $\varkappa_1$ theilbar und folglich mit $\varkappa_1$ associirt sein. Dann sind auch die Producte
\[
\pi_2 \ldots \pi_\nu, \qquad \varkappa_2 \ldots \varkappa_\mu
\]
associirt, und folglich ist einer der Factoren des ersten Productes, etwa $\pi_2$, durch $\varkappa_2$ theilbar und daher mit $\varkappa_2$ associirt. So kann man weiter schliessen, und es ergiebt sich, dass nicht nur die Anzahl der $\varkappa$ mit der Anzahl der $\pi$ übereinstimmen muss, sondern dass auch die $\varkappa$ einzeln den $\pi$ associirt sind.

Unter den Primfactoren eines ganzen Functionales $\omega$ kann derselbe mehrmals vorkommen, und diese einander gleichen (oder associirten) Factoren können zu einer Potenz zusammengefasst werden. Ist $\omega$ durch $\pi^h$ theilbar, so sagen wir, \gesperrt{das Primfunctional $\pi$ geht $h$ mal in $\omega$ auf.}

Hiernach hat es einen ganz bestimmten Sinn, wenn von \gesperrt{den Primfactoren} eines ganzen Functionales gesprochen wird. Wir folgern noch aus dem Bewiesenen:

4. \gesperrt{Ein ganzes Functional $\alpha$ ist dann und nur dann durch ein anderes $\beta$ theilbar, wenn alle Primfactoren von $\beta$ unter den Primfactoren von $\alpha$ vorkommen, und jeder von ihnen mindestens so oft in $\alpha$ aufgeht, als in $\beta$.}

Denn ist $\alpha$ durch $\beta$ und $\beta$ durch $\pi^h$ theilbar, so ist auch $\alpha$ durch $\pi^h$ theilbar.

Sind $\alpha, \beta, \ldots$ ganze Functionale, $\pi_1, \pi_2, \ldots$ verschiedene Primfunctionale, $\varepsilon_1, \varepsilon_2, \ldots$ Einheiten, so können wir Exponenten $a_1, a_2, \ldots, b_1, b_2, \ldots$ so bestimmen, dass
\weq{7}{\begin{aligned}
\varepsilon_1 \alpha &= \pi_1^{a_1} \pi_2^{a_2} \ldots \\
\varepsilon_2 \beta &= \pi_1^{b_1} \pi_2^{b_2} \ldots \\
&\ldots\ldots\ldots\ldots
\end{aligned}}
%CONT>>
